% =====================================================================
%  Unsupervised Test-Time Adaptation for Cross-Corpus (Multilingual)
%  Audio Deepfake Detection. Compiles on Overleaf as-is (IEEEtran).
% =====================================================================
\documentclass[conference]{IEEEtran}

\usepackage{cite}
\usepackage{amsmath}
\usepackage{amssymb}
\usepackage{algorithm}
\usepackage{algpseudocode}
\usepackage{graphicx}
\usepackage{booktabs}
\usepackage{tabularx}
\usepackage[table]{xcolor}
\usepackage{textcomp}

% include a figure if present, else a placeholder box (compiles before upload)
\newcommand{\safefig}[2]{%
  \IfFileExists{#2}{\includegraphics[width=#1]{#2}}{%
    \fbox{\parbox[c][2.6cm][c]{\dimexpr#1-2\fboxsep-2\fboxrule\relax}{%
      \centering\ttfamily\detokenize{#2}\\[3pt]\normalfont(upload this figure)}}}%
}

\begin{document}

\title{Ranking Transfers, Thresholds Don't:\\
Unsupervised Test-Time Adaptation for Cross-Corpus Audio Deepfake Detection}

\author{\IEEEauthorblockN{Roham Izadidoost}
\IEEEauthorblockA{Manipal University Jaipur \\
izadidoostroham@gmail.com}
\and
\IEEEauthorblockN{Dr.~Sumit Srivastava}
\IEEEauthorblockA{Manipal University Jaipur}
\and
\IEEEauthorblockN{Shweta Sharma}
\IEEEauthorblockA{Manipal University Jaipur \\
Shweta.sharma1@jaipur.manipal.edu}}

\maketitle

\begin{abstract}
Audio deepfake detectors report near-perfect in-domain accuracy but degrade
sharply on speech from unseen corpora---different recording conditions,
generators, and languages---limiting real-world use. We make a simple empirical
observation with large consequences: across corpora, a fine-tuned
self-supervised detector's \emph{ranking} (ROC-AUC) largely transfers even when
its decision \emph{threshold} does not calibrate. This means the model's own
most-confident predictions on unlabeled target audio are reliable, which in turn
enables stable unsupervised \emph{test-time adaptation} (TTA): we self-train on
confident-ranked pseudo-labels and enforce channel-consistency, adapting only
LayerNorm and the classifier head---no target labels and no source data at test
time. On held-out In-the-Wild the method cuts Equal Error Rate (EER) from
9.4\% to 6.4\% (stable over three seeds), with a smaller gain on ASVspoof2021-DF
where the domain gap is smaller, and it generalises inductively to unseen target
clips. In contrast, entropy-minimisation TTA collapses to chance, and train-time
domain-adversarial objectives do not help. We further extend the study to Arabic
for a multilingual cross-corpus evaluation.
\end{abstract}

\begin{IEEEkeywords}
audio deepfake detection, anti-spoofing, test-time adaptation, cross-corpus
generalisation, self-supervised learning, EER.
\end{IEEEkeywords}

\section{Introduction}
Synthetic speech is realistic enough to threaten voice authentication and enable
fraud, motivating reliable spoof/deepfake detectors. On standard benchmarks,
self-supervised (SSL) front-ends such as wav2vec\,2.0 and XLS-R, paired with
lightweight back-ends, reach very low error~\cite{wav2vec2aug,xlsrsls,aasist}.
Yet these numbers are largely \emph{in-domain}: trained and tested on the same
corpus. Deployed on audio from a different source---new microphones, codecs,
languages, or generators---accuracy drops steeply, a well-documented
generalisation gap~\cite{asdg,crossdomain}.

We start from a diagnostic measurement. Taking a fine-tuned XLS-R detector and
evaluating it on an unseen corpus, we observe that its \textbf{ROC-AUC stays
high while its accuracy at a fixed threshold collapses}: the model still
\emph{ranks} clips correctly (real below fake) but its operating point is
mis-calibrated for the new domain. This asymmetry is the key that unlocks a
practical fix. Because the ranking transfers, the model's most-confident
predictions on unlabeled target audio are trustworthy, so we can use them to
adapt the model to the target domain \emph{at test time}, without any labels.

Our contributions are:
\begin{itemize}
\item An empirical finding---\emph{ranking transfers cross-corpus even when the
threshold does not}---and its consequence for unsupervised adaptation.
\item A simple, stable TTA method: confident-ranked pseudo-label self-training
plus channel-consistency, adapting only LayerNorm and the head.
\item A cross-corpus, multilingual evaluation (In-the-Wild, ASVspoof2021-DF,
Arabic) with multi-seed, ablation and inductive analyses, and a clear contrast
against entropy-minimisation TTA (which collapses) and train-time
domain-adversarial training (which does not help).
\end{itemize}

\section{Related Work}
\textbf{SSL anti-spoofing.} Frozen or fine-tuned wav2vec\,2.0 / XLS-R / WavLM
front-ends with AASIST or attentive back-ends dominate recent
benchmarks~\cite{wav2vec2aug,xlsrsls,aasist,rawnet2}. These are strong in-domain
but corpus-specific.
\textbf{Domain generalisation (DG) for ADD.} Train-time methods align real and
separate fake across source domains (e.g.\ aggregation-separation~\cite{asdg}) or
fuse SSL features~\cite{crossdomain}. They target the gap at training time and
cannot use information about the deployment domain.
\textbf{Test-time adaptation.} TTA adapts a trained model to an unlabeled target
distribution, e.g.\ entropy minimisation over normalisation
parameters~\cite{tent}. It is widely used in vision but under-explored for audio
deepfake detection; we show naive entropy minimisation is unstable here and
propose a stable, ranking-driven alternative.

\section{Method}
\subsection{Problem formulation}
Let $x$ be a waveform with label $y\in\{0,1\}$ ($0$ real, $1$ fake). A detector
$f_\theta$ maps $x$ to frame features, pools them to an embedding
$e\in\mathbb{R}^{d}$, and outputs a fake probability
\begin{equation}
s_\theta(x)=\sigma\!\left(w^\top e+b\right)=P_\theta(y{=}1\mid x),\qquad \sigma(z)=\tfrac{1}{1+e^{-z}} .
\end{equation}
Training on a labeled \emph{source} corpus $\mathcal{D}_S=\{(x_i,y_i)\}$ minimises
the cross-entropy
\begin{equation}
\mathcal{L}_{\mathrm{CE}}(\theta;\mathcal{D})=-\frac{1}{|\mathcal{D}|}\!\!\sum_{(x,y)\in\mathcal{D}}\!\!
\big[y\log s_\theta(x)+(1{-}y)\log(1{-}s_\theta(x))\big].
\end{equation}
At test time we receive an \emph{unlabeled} target corpus
$\mathcal{D}_T=\{x_j\}_{j=1}^{N}$ from a shifted distribution (different corpus,
channel, or language); no target labels and no source data are available.

\subsection{Detector}
The encoder $g_\psi$ is XLS-R-300M with only its top four transformer layers
trainable. Given frame features $\{h_t\}_{t=1}^{T}$, attentive temporal pooling
forms the embedding
\begin{equation}
\alpha_t=\frac{\exp(v^\top h_t)}{\sum_{t'}\exp(v^\top h_{t'})},\qquad
e=\sum_{t=1}^{T}\alpha_t\,h_t,
\end{equation}
followed by a projection and the linear head of Eq.~(1). Fine-tuning only the top
layers closes most of the in-domain gap while remaining trainable on one GPU.

\subsection{The ranking--calibration gap}
For a threshold $\tau$ the decision is $\hat y_\tau(x)=\mathbf{1}[s_\theta(x)\ge\tau]$,
whose accuracy depends on $\tau$. Ranking quality, by contrast, is threshold-free:
\begin{equation}
\mathrm{AUC}=\Pr_{x^{+}\sim\mathrm{fake},\,x^{-}\sim\mathrm{real}}\!\big[s_\theta(x^{+})>s_\theta(x^{-})\big],
\end{equation}
which is invariant to any strictly increasing reparametrisation of $s_\theta$.
Empirically, under the source$\rightarrow$target shift $\mathrm{AUC}$ is largely
preserved while the accuracy-optimal threshold $\tau^\star$ drifts far from $0.5$
(Fig.~\ref{fig:scores}: AUC $0.966$ yet accuracy at $\tau{=}0.5$ is $61\%$).
\emph{Ranking transfers; calibration does not.} We therefore trust the score
\emph{order}, not its absolute value: the extreme quantiles are high-precision
pseudo-labels even when $s_\theta$ is mis-calibrated.

\begin{figure*}[t]
\centering
\safefig{0.86\textwidth}{fig_score_dist.png}
\caption{Detector scores (logit axis; scores concentrate near~1) for real vs.\
fake clips on the \emph{unseen} In-the-Wild corpus. \textbf{Left:} the source
model pushes almost all clips past the $0.5$ threshold---78\% of \emph{real}
clips are scored as fake (acc@0.5$=$61\%)---yet real still ranks below fake
(AUC$=$0.966): ranking transfers, calibration does not. \textbf{Right:}
unsupervised test-time adaptation shifts the real mass back across the threshold
(acc@0.5$=$93\%, AUC$=$0.987). Representative seed; Table~\ref{tab:main} reports
the 3-seed mean EER.}
\label{fig:scores}
\end{figure*}

\subsection{Test-time adaptation}
Let $Q_\alpha$ denote the $\alpha$-quantile of the target scores
$\{s_\theta(x_j)\}_{j=1}^{N}$. Each adaptation epoch we (re-)assign confident
pseudo-labels from the extreme tails,
\begin{equation}
\tilde y_j=
\begin{cases}
0, & s_\theta(x_j)\le Q_{q},\\
1, & s_\theta(x_j)\ge Q_{1-q},\\
\varnothing, & \text{otherwise (ignored)},
\end{cases}
\end{equation}
and let $\mathcal{C}=\{j:\tilde y_j\neq\varnothing\}$. Self-training is
cross-entropy on this confident set,
\begin{equation}
\mathcal{L}_{\mathrm{ST}}=-\frac{1}{|\mathcal{C}|}\sum_{j\in\mathcal{C}}
\big[\tilde y_j\log s_\theta(x_j)+(1{-}\tilde y_j)\log(1{-}s_\theta(x_j))\big].
\end{equation}
With a stochastic channel augmentation $a(\cdot)$ (random gain $+$ additive noise)
and $p_\theta(x)=[\,1{-}s_\theta(x),\,s_\theta(x)\,]$, a consistency term enforces
prediction invariance to the channel,
\begin{equation}
\mathcal{L}_{\mathrm{cons}}=\frac{1}{N}\sum_{j=1}^{N}\big\|\,p_\theta(a(x_j))-\mathrm{sg}\!\left[p_\theta(x_j)\right]\big\|_2^2,
\end{equation}
where $\mathrm{sg}[\cdot]$ is stop-gradient. The adaptation objective is
\begin{equation}
\min_{\{\gamma,\beta\},\,\phi}\ \ \mathcal{L}_{\mathrm{ST}}+\lambda\,\mathcal{L}_{\mathrm{cons}},
\qquad \lambda=0.3,\ q=0.3,
\end{equation}
optimising \emph{only} the LayerNorm affine parameters $\{\gamma,\beta\}$ of the
fine-tuned layers and the head $\phi$; the rest of $\psi$ stays frozen. Restricting
the trainable set and using only confident tails prevents the degenerate collapse
$s_\theta(\cdot)\!\to\!\text{const}$ to which entropy minimisation is prone, as we
confirm empirically against Tent (Sec.~V). Algorithm~\ref{alg:tta} summarises the
procedure.

\begin{algorithm}[t]
\caption{Real-manifold test-time adaptation (per target corpus)}
\label{alg:tta}
\begin{algorithmic}[1]
\Require source model $f_\theta$; unlabeled target $\mathcal{D}_T=\{x_j\}_{j=1}^{N}$;
quantile $q$; weight $\lambda$; epochs $E$; augmentation $a(\cdot)$; step size $\eta$
\Ensure adapted model $f_\theta$
\State make top-layer LayerNorm $\{\gamma,\beta\}$ and head $\phi$ trainable; freeze all else
\For{$e=1,\dots,E$}
  \State $s_j \gets s_\theta(x_j)\ \ \forall j$ \Comment{score the whole target set}
  \State $Q_q,\,Q_{1-q}\gets$ quantiles of $\{s_j\}_{j=1}^{N}$
  \State $\tilde y_j \gets 0$ if $s_j\!\le\!Q_q$;\ \ $1$ if $s_j\!\ge\!Q_{1-q}$;\ \ else $\varnothing$
  \State $\mathcal{C}\gets\{j:\tilde y_j\neq\varnothing\}$ \Comment{confident tails}
  \For{minibatch $\mathcal{B}\subset\mathcal{D}_T$}
    \State $\mathcal{L}_{\mathrm{ST}}\gets\mathrm{CE}$ over $\mathcal{B}\cap\mathcal{C}$ with targets $\tilde y$ \Comment{Eq.~(6)}
    \State $\mathcal{L}_{\mathrm{cons}}\gets\frac{1}{|\mathcal{B}|}\sum_{x\in\mathcal{B}}\|p_\theta(a(x))-\mathrm{sg}[p_\theta(x)]\|_2^2$
    \State $\theta\gets\theta-\eta\,\nabla_\theta(\mathcal{L}_{\mathrm{ST}}+\lambda\mathcal{L}_{\mathrm{cons}})$
  \EndFor
\EndFor
\end{algorithmic}
\end{algorithm}

\section{Experimental Setup}
\textbf{Source.} ASVspoof2019-LA (train), class-balanced.
\textbf{Targets (unseen).} In-the-Wild (English, real+fake), ASVspoof2021-DF, and
the Arabic Audio Deepfake set (ArAD) for a multilingual target.
\textbf{Metric.} Equal Error Rate (EER) and ROC-AUC. With false-positive and
false-negative rates $\mathrm{FPR}(\tau),\mathrm{FNR}(\tau)$,
\begin{equation}
\mathrm{EER}=\mathrm{FPR}(\tau^\star)=\mathrm{FNR}(\tau^\star),\quad
\tau^\star:\ \mathrm{FPR}(\tau^\star)=\mathrm{FNR}(\tau^\star),
\end{equation}
which is threshold-free and thus insensitive to the calibration drift above.
\textbf{Setting.} Transductive TTA (adapt on the unlabeled target pool, then
score it); we additionally report an inductive check (adapt on one split, test on
a disjoint one). \textbf{Implementation.} 16\,kHz, 3\,s crops; bf16; Adam,
lr $10^{-4}$; top-4-layer fine-tuning; single GPU.

\section{Results}
\subsection{Cross-corpus adaptation}
Table~\ref{tab:main} is the main result. On In-the-Wild, where the domain gap is
large, TTA reduces EER from 9.42\% to 6.41\% ($\pm0.11$ over three seeds) and
raises AUC. On ASVspoof2021-DF, closer to the source distribution, the source is
already strong (5.50\%) and the gain is small (5.29\%)---the method helps in
proportion to the gap. Vanilla entropy-minimisation TTA (Tent) collapses to
chance, underlining the importance of the confident-pseudo-label design.

\begin{table}[t]
\centering
\caption{Cross-corpus EER\,\%\,/\,AUC (lower EER better). Ours = confident
self-training + consistency (mean over 3 seeds). Arabic to be completed from the
full multilingual run.}
\label{tab:main}
\begin{tabular}{@{}lcc@{}}
\toprule
Method & In-the-Wild & ASVspoof2021-DF \\
\midrule
Source-only (no adaptation) & 9.42 / 0.966 & 5.50 / 0.989 \\
Tent (entropy-min TTA)~\cite{tent} & 49.45 / 0.506 & --- \\
\rowcolor{green!12}\textbf{Ours (TTA)} & \textbf{6.41 / 0.985} & \textbf{5.29 / 0.989} \\
\midrule
\multicolumn{3}{@{}l}{\emph{Arabic (ArAD), multilingual target:} source X.XX / ours Y.YY (pending)}\\
\bottomrule
\end{tabular}
\end{table}

\subsection{Ablation}
Table~\ref{tab:abl} isolates each component on In-the-Wild. Confident
self-training alone gives most of the gain; channel-consistency adds a further
improvement and yields the best configuration. A real-anchoring contrastive term
(pulling target-real embeddings to a source real prototype) does \emph{not} help
and is dropped.

\begin{table}[t]
\centering
\caption{Ablation on In-the-Wild (EER\,\%\,/\,AUC, seed 0).}
\label{tab:abl}
\begin{tabular}{@{}lc@{}}
\toprule
Configuration & In-the-Wild \\
\midrule
Source-only & 9.42 / 0.966 \\
Self-training only & 7.13 / 0.983 \\
\quad + real-anchor & 7.37 / 0.982 \\
\rowcolor{green!12}\quad + consistency (\textbf{ours}) & \textbf{6.37 / 0.985} \\
Self-train + anchor + consistency & 6.80 / 0.984 \\
\bottomrule
\end{tabular}
\end{table}

\subsection{Inductive check}
To rule out transductive memorisation, we adapt on one In-the-Wild split and
evaluate on a \emph{disjoint} split: EER improves 9.88\%$\rightarrow$7.48\%. The
adaptation therefore captures domain-level structure that transfers to unseen
target clips, not just the specific evaluation items.

\subsection{What did not work}
Train-time domain-adversarial training (gradient reversal on corpus identity)
plus a real-anchored contrastive loss---our initial approach---did not improve
cross-corpus EER and was unstable (adversarial training collapsed). This negative
result motivated the shift to test-time adaptation and is reported for
transparency.

\section{Discussion and Limitations}
The method is unsupervised in labels but \emph{transductive}: it uses unlabeled
target inputs at adaptation time, a standard TTA setting that we state explicitly;
the inductive check addresses the memorisation concern. Gains are largest under
large domain shift and negligible when the source already matches the target,
which is expected. The approach is cheap (adapt a trained model in minutes) and
orthogonal to the choice of SSL back-end.

\section{Conclusion}
Cross-corpus audio deepfake detection fails not because ranking is lost but
because the operating point is mis-calibrated. Exploiting this, a simple
unsupervised test-time adaptation---confident-ranked self-training plus
channel-consistency---reliably improves generalisation to unseen corpora and
languages, where entropy-minimisation TTA collapses and train-time
domain-adversarial training does not help.

\begin{thebibliography}{1}
\bibitem{wav2vec2aug} H.~Tak \emph{et al.}, ``Automatic speaker verification
spoofing and deepfake detection using wav2vec 2.0 and data augmentation,'' in
\emph{Proc.\ Odyssey}, 2022.
\bibitem{xlsrsls} Q.~Zhang \emph{et al.}, ``Audio deepfake detection with
self-supervised XLS-R and SLS classifier,'' 2024.
\bibitem{aasist} J.~Jung \emph{et al.}, ``AASIST: Audio anti-spoofing using
integrated spectro-temporal graph attention networks,'' in \emph{ICASSP}, 2022.
\bibitem{rawnet2} H.~Tak \emph{et al.}, ``End-to-end anti-spoofing with RawNet2,''
in \emph{ICASSP}, 2021.
\bibitem{asdg} Y.~Xie \emph{et al.}, ``Domain generalization via aggregation and
separation for audio deepfake detection,'' \emph{IEEE TIFS}, 2024.
\bibitem{crossdomain} ``Cross-domain audio deepfake detection: dataset and
analysis,'' 2024.
\bibitem{tent} D.~Wang \emph{et al.}, ``Tent: Fully test-time adaptation by
entropy minimization,'' in \emph{ICLR}, 2021.
\end{thebibliography}

\end{document}
